% Evaluation

This section evaluates the \emph{technical and economic feasibility} of AnimatorSBT
along three dimensions:
marginal-cost analysis through cost simulation (\cref{subsec:cost_sim}),
scalability analysis (\cref{subsec:scalability}),
and a qualitative assessment against the requirements defined in \cref{subsec:formalization}.
We emphasize at the outset that this is a feasibility-and-correctness evaluation,
not a validation of real-world adoption or usability: the prototype shows that the
mechanism works and is inexpensive to operate, but claims about animator uptake,
perceived value, and studio participation remain to be tested
(\cref{subsec:eval_limitations}).

% ----- 6.1 Cost Simulation -----
\subsection{Cost Simulation}
\label{subsec:cost_sim}

We simulated the annual operating cost of AnimatorSBT across five adoption scenarios,
ranging from a single-studio pilot to full industry adoption.
The simulation uses industry statistics from JAniCA~(2023) and NAFCA~(2024)
for workforce parameters, and gas costs measured on-chain on the Polygon Amoy
testnet (\cref{tab:gas}; 30~gwei gas price, POL = \$0.22, March 2026 rate).

\subsubsection{Parameters}

The simulation assumes the following industry parameters based on public data:

\begin{itemize}
  \item Total animation creators: $\sim$5{,}500 (JAniCA 2023 survey base)
  \item Freelance ratio: 47.3\% (JAniCA 2023) to 69.6\% (JAniCA 2019). The
    high-estimate scenario applies the older 2019 ratio to the 2023 creator base
    as an upper bound; this mixes survey years and should be read as a conservative
    ceiling, not a current figure.
  \item Average projects per animator per year: 6
  \item Batch minting with batch size of 10 (default)
  \item Revocation rate: 1\% of issued SBTs
  \item IPFS pinning: Free tier for $\leq$500 files; \$20/month for larger deployments~\cite{pinata2026pricing}
\end{itemize}

\subsubsection{Results}

\Cref{tab:cost_scenarios} presents the simulation results.

\begin{table}[ht]
\centering
\caption{Annual cost simulation across five adoption scenarios.
All costs in USD at Polygon PoS rates (30 gwei, POL = \$0.22, March 2026).
Gas values are the on-chain Amoy measurements (\cref{tab:gas}); USD figures are
projections (see \cref{subsec:eval_limitations}).}
\label{tab:cost_scenarios}
\footnotesize
\setlength{\tabcolsep}{3pt}
\begin{tabular}{@{}lrrrr@{}}
\toprule
\textbf{Scenario} & \textbf{Anim.} & \textbf{SBTs/yr} & \textbf{Cost/yr} & \textbf{Per Anim.} \\
\midrule
Pilot (1 studio)      & 50    & 300      & \$0.22  & \$0.004 \\
Early (10 studios)    & 500   & 3{,}000  & \$242   & \$0.484 \\
Moderate (50\%)       & 1{,}300 & 7{,}800 & \$245  & \$0.189 \\
Full (low est.)       & 2{,}601 & 15{,}606 & \$251 & \$0.096 \\
Full (high est.)      & 3{,}827 & 22{,}962 & \$256 & \$0.067 \\
\bottomrule
\end{tabular}
\end{table}

\Cref{fig:cost_simulation} visualizes the cost breakdown and economies of scale.

\begin{figure}[t]
  \centering
  \includegraphics[width=\columnwidth]{figures/cost_simulation}
  \caption{Cost simulation results. Left: annual cost breakdown by component
  across five adoption scenarios. Right: per-animator cost showing
  economies of scale as adoption increases.}
  \label{fig:cost_simulation}
\end{figure}

The key finding is that the \emph{marginal} operating cost is low and---importantly---is
\textbf{dominated by a flat off-chain SaaS fee, not by the blockchain}.
At full industry adoption the modeled total is approximately \$256/year, of which
roughly \$240 is the fixed IPFS pinning subscription~\cite{pinata2026pricing}
and only about \$16 is on-chain minting.
The jump from the Pilot scenario (\$0.22) to the Early scenario (\$242) is caused
entirely by crossing Pinata's free-tier threshold (300 vs.\ 3{,}000 files),
not by blockchain costs. The blockchain layer itself is inexpensive: Layer~2 fees
make per-transaction cost economically negligible.

We are careful not to overstate what this figure means.
First, it is a \emph{marginal} cost (on-chain transactions plus IPFS storage);
it excludes the operating costs the system actually requires at scale---running the
Issuance Service, \textbf{a relayer that sponsors gas on behalf of all animators}
(\cref{subsec:flow}), wallet-recovery support, the co-signing workflow
(\cref{subsec:cosigning}), PMC integration and maintenance, and the legal/DPIA
review we recommend (\cref{subsec:legal}).
Second, because the relayer subsidizes gas, the \$256 is borne by a single
sponsoring party fronting costs for the entire industry on an ongoing basis---a
funding-sustainability and centralization consideration, not a free lunch.
We therefore present \$256 as evidence that \emph{the blockchain layer is not a
cost barrier}, not as a total cost of ownership.
We also avoid normalizing the figure against industry revenue, since that revenue
accrues largely to production committees and IP holders rather than to the freelance
animators the system serves; the relevant comparison is the operating budget of
whoever runs the service.

The per-animator marginal cost decreases with scale as the fixed IPFS cost is
spread across more users, reaching roughly \$0.067 per animator per year at full
adoption.

Because the dollar figures depend on two volatile parameters---gas price and the
POL/USD rate---\Cref{tab:sensitivity} reports the per-mint cost over a range of
each, computed arithmetically from the on-chain mint gas (189{,}792). The
parameter values are \emph{hypothetical scenarios}, not observed market prices.
Even in the most adverse cell (300~gwei, POL at \$1.00), a single mint costs under
\$0.06; the qualitative conclusion---that on-chain cost is not a barrier---is robust
to large swings in both parameters.

\begin{table}[ht]
\centering
\caption{Sensitivity of the per-\texttt{mint()} cost (USD) to gas price and POL/USD
rate, computed from the on-chain mint gas of 189{,}792. Values are hypothetical
price scenarios, not observed market rates. The shaded baseline used elsewhere in
this paper is 30~gwei / \$0.22.}
\label{tab:sensitivity}
\footnotesize
\setlength{\tabcolsep}{6pt}
\begin{tabular}{@{}lrrr@{}}
\toprule
\textbf{Gas price} & \textbf{POL \$0.22} & \textbf{POL \$0.50} & \textbf{POL \$1.00} \\
\midrule
30~gwei  & \$0.0013 & \$0.0028 & \$0.0057 \\
100~gwei & \$0.0042 & \$0.0095 & \$0.0190 \\
300~gwei & \$0.0125 & \$0.0285 & \$0.0569 \\
\bottomrule
\end{tabular}
\end{table}

% ----- 6.2 Scalability -----
\subsection{Scalability Analysis}
\label{subsec:scalability}

\subsubsection{Transaction Throughput}

Polygon PoS processes approximately 3--8 million transactions per day
with a block time of approximately 2 seconds~\cite{gangwal2023layer2}.
At full adoption, AnimatorSBT would generate approximately 22{,}962 minting
transactions per year, or roughly 63 per day.
This represents less than 0.002\% of Polygon's daily capacity,
indicating that the system operates well within the network's limits.

\subsubsection{Storage}

Each SBT metadata JSON file is approximately 500 bytes.
At full adoption, annual metadata storage is approximately
$22{,}962 \times 500 = 11.5$~MB---trivially small for IPFS.
On-chain storage is minimal (token ID, owner address, URI string),
at approximately 200 bytes per token.

\subsubsection{On-Chain Query Cost}

The \texttt{tokensOfOwner()} function performs a linear scan
across all minted tokens ($O(n)$ where $n$ is the total token count).
While acceptable at the current scale (tens of thousands of tokens),
this becomes inefficient as $n$ grows:
at 100{,}000 tokens, the function may exceed RPC node timeout limits
for \texttt{eth\_call} queries.
A production deployment should add an owner-to-token-ID mapping
for $O(1)$ lookups, or rely on off-chain indexing
(e.g., The Graph) for portfolio queries.
Crucially, it is the \emph{cumulative} token count, not the annual flow, that
drives this limit: at the projected $22{,}000$--$42{,}000$ SBTs per year
(\cref{subsec:scalability}), cumulative supply enters the $\sim$100k range within a
few years of deployment. We therefore treat the prototype's \texttt{tokensOfOwner()}
as demonstration-only and regard the enumerable mapping or an off-chain indexer as a
launch requirement rather than a later optimization. (We do not report gas for the
mapping variant here, as it would require re-measurement of the modified contract.)

\subsubsection{Growth Projection}

Assuming a 10.6\% CAGR in the anime market~\cite{mordor2026anime}
and proportional workforce growth,
the system would need to handle approximately 42{,}000 SBTs per year by 2031
(roughly 115 per day).
This remains well within Polygon's capacity
and would cost approximately \$300--400 annually---again dominated by the flat
IPFS pinning tier rather than on-chain fees, and subject to the same
relayer-subsidy and total-cost-of-ownership caveats noted for the \$256 figure
(\cref{subsec:cost_sim}).

% ----- 6.3 Requirements Evaluation -----
\subsection{Requirements Satisfaction}
\label{subsec:requirements_eval}

We assess AnimatorSBT against the five requirements
defined in \cref{subsec:formalization}.
We distinguish what the prototype \emph{empirically demonstrates}---the
integrity/existence aspect of R2 and the cost and throughput results of
\cref{subsec:cost_sim,subsec:scalability}---from what it satisfies \emph{by design
but has not yet validated} with users, at scale, or with the unbuilt components
(R1, R4, R5, and the authenticity aspect of R2).

\begin{itemize}
  \item \textbf{R1 (Completeness):} \textit{Satisfied by design, with a noted tension.}
    Any animator using the PMC can receive an SBT regardless of their
    position in the subcontracting hierarchy, because issuance is triggered by
    work-log data in the PMC, not by studio approval or credit-committee decisions.
    Two caveats: (i)~completeness presupposes PMC adoption---an animator who does
    not use the tool remains invisible; and (ii)~the open, low-friction
    self-enrollment that R1 requires is in tension with the Sybil mitigation of
    \cref{subsec:trust}, which treats PMC registration as an identity gate. One
    cannot simultaneously have fully open self-onboarding and PMC registration as a
    meaningful Sybil barrier; we revisit this trade-off in \cref{subsec:trust}.

  \item \textbf{R2 (Verifiability):} \textit{Partially satisfied (integrity and
    existence, not authenticity).}
    Third parties can verify, by querying the public blockchain and recomputing the
    evidence hash, that a contribution claim \emph{exists}, is \emph{unaltered}, and
    carries a \emph{timestamp}, without contacting the original studio. They cannot,
    from the chain alone, verify that the claimed work was actually performed: the
    metadata derives from a self-reported log, and \texttt{evidence\_hash} is
    computed over issuer-authored attributes (\cref{subsec:verification,subsec:trust}).
    Authenticity requires the studio co-signing layer of \cref{subsec:cosigning},
    which we now provide as a tested reference contract but have not yet deployed
    on-chain or operated with an \emph{independent, identity-bound} studio attester.
    Until that independent operation exists, a co-signed record is
    \emph{operator-attested}, not independently authenticated. This gap---from a
    working mechanism to trustworthy independent attestation---is the single most
    important qualification in this paper.

  \item \textbf{R3 (Persistence):} \textit{Partially satisfied.}
    On-chain data (token ownership, issuance timestamp) persists
    as long as the Polygon network operates.
    However, the detailed metadata stored on IPFS depends on
    continued pinning by at least one party.
    If the pinning service subscription lapses and no other node
    holds a copy, the metadata becomes inaccessible
    despite the on-chain record remaining intact.
    This is a known limitation of IPFS-based storage;
    mitigation strategies include multi-provider pinning,
    community-operated IPFS nodes, and Arweave-based
    permanent storage as a fallback.

  \item \textbf{R4 (Portability):} \textit{Satisfied by design, pending implementation.}
    The SBT is bound to the animator's wallet address, which they control
    via account abstraction, and can be presented to any verifier independent of
    any platform. This rests, however, on the embedded-wallet layer that is not yet
    built (\cref{subsec:stack}), and the unsolved wallet-loss/recovery problem
    (\cref{subsec:trust}) directly threatens portability: losing wallet access today
    would mean losing the entire portfolio. We thus regard R4 as satisfied in design
    but contingent on a recovery mechanism that remains future work.

  \item \textbf{R5 (Minimal Friction):} \textit{Satisfied by design, not yet validated.}
    In the intended design the animator's workflow is unchanged and issuance is
    automatic upon project milestones. Three qualifications apply: (i)~the
    walletless, gasless experience depends on ERC-4337 and relayer components that
    are unimplemented (\cref{subsec:stack}); (ii)~the system also \emph{introduces}
    friction not captured by ``unchanged workflow''---wallet provisioning, issuance
    consent, and managing selective disclosure; and (iii)~no user study has yet
    measured perceived friction (\cref{subsec:eval_limitations}). We therefore mark
    R5 as satisfied by design but requiring empirical user validation.
\end{itemize}

% ----- 6.4 Limitations -----
\subsection{Limitations}
\label{subsec:eval_limitations}

The current evaluation has several limitations that should be addressed
in future work:

\begin{itemize}
  \item \textbf{Testnet deployment only:}
    The prototype has been deployed on the Polygon Amoy testnet,
    not on mainnet.
    While gas costs are measured accurately (the EVM execution
    is identical), a mainnet deployment would be required
    for production use.
  \item \textbf{No user study:}
    We have not yet conducted a formal user study with practicing animators.
    The cost simulation and measured gas data provide quantitative validation,
    but real-world usability and perceived value remain to be evaluated.
  \item \textbf{Cost figures are testnet-derived projections:}
    We measured gas \emph{used} on the Polygon Amoy testnet, which transfers
    faithfully to mainnet because EVM execution is identical. The USD figures,
    however, are \emph{projections}: they apply a fixed gas price (30~gwei) and a
    fixed POL/USD rate (\$0.22) to that gas, and they exclude mainnet priority fees
    and any L1 data costs. Testnet POL has no market value, so the dollar column is
    a model, not an observed cost. We provide a sensitivity analysis over gas and
    POL prices (\cref{tab:sensitivity}); the qualitative conclusion is robust, but
    the specific dollar figures should be read as projections rather than measured
    market costs.
  \item \textbf{Self-reported data and a self-referential evidence hash:}
    The SBT metadata is derived from self-reported work logs in the PMC.
    We have now implemented the co-signing mechanism as a reference contract
    (\cref{subsec:cosigning}), but until it is deployed and operated by an
    independent, identity-bound studio it adds no authenticity, so as evaluated the
    system still cannot independently verify the truth of contribution claims---only
    their existence and integrity. The limitation is in fact sharper than ``no
    co-signing'': \texttt{evidence\_hash} is a hash of the issuer-authored attributes
    stored in the \emph{same} metadata object, so recomputing it from that object
    detects only post-mint tampering, not correspondence to any independent source
    (\cref{subsec:verification}). Comparison against the animator's local work log
    is stronger but still self-referential to a self-reported log. The
    serialization/hashing pipeline that this comparison relies on \emph{is} now
    tested (7 canonical-serialization unit tests, \cref{subsec:prototype}), closing a
    determinism gap noted earlier; what remains untested is the live IPFS path and
    the in-browser portal recompute. We discuss the trust model and mitigation
    strategies in \cref{sec:discussion}.

  \item \textbf{Wallet loss and recovery (unsolved):}
    Because SBTs are bound to a wallet address, loss of wallet access means loss of
    all accumulated credentials. The social-recovery / \texttt{migrate()} direction
    discussed in \cref{subsec:trust} is \emph{unimplemented} and is an open problem
    in the broader SBT literature. This directly undercuts R3 (persistence of an
    \emph{accessible} record) and R4 (portability), and is a prerequisite for any
    real deployment.
  \item \textbf{IPFS metadata persistence:}
    Detailed SBT metadata resides on IPFS, which requires active pinning
    to remain accessible. If all pinning providers stop hosting the data,
    the on-chain token record persists but the contribution details
    become unrecoverable.
  \item \textbf{Adoption assumptions:}
    The simulation assumes uniform adoption across the industry.
    In practice, adoption will be uneven and influenced by
    network effects, studio policies, and cultural factors.
\end{itemize}
